\documentclass[11pt,a4paper]{article}
\usepackage[utf8]{inputenc}
\usepackage[T1]{fontenc}
\usepackage{authblk}
\usepackage{hyperref}
\usepackage{graphicx}
\usepackage{amsmath}
\usepackage{amssymb}
\usepackage{geometry}
\usepackage{listings}
\usepackage{xcolor}
\usepackage{booktabs}
\usepackage{abstract}

\geometry{margin=1in}

\hypersetup{
    colorlinks=true,
    linkcolor=blue,
    filecolor=magenta,      
    urlcolor=cyan,
    pdftitle={Aura: Sovereign AI Research (aura.meziani.org)},
    pdfpagemode=FullScreen,
}

\title{Aura: A Sovereign Agentic Monorepo for \\ Safe AI Research and Deployment}

\author[1]{Yani Meziani \thanks{ORCID: 0009-0007-4348-8711}}
\affil[1]{Aura Research Lab (aura.meziani.org)}

\date{\today}

\begin{document}

\maketitle

\begin{abstract}
This manifesto outlines the architecture and deployment strategy for Aura, an open-source, sovereign, and military-grade decentralized agentic monorepo. Building upon foundational work in multimodal architectures and Hamiltonian State Space Duality (as seen in Akasha 2 \cite{meziani2026akasha}), Aura integrates mobile control (Pegasus), cloud execution (Cerberus), and specialized AI agents into a unified, zero-trust ecosystem. We prioritize operator-first design and extreme resilience to ensure AI remains a tool for human empowerment.
\end{abstract}

\section{Introduction}
Aura is a sovereign system deployment that emphasizes local-first control and military-grade security protocols. By integrating a lean Zig-based agent runtime (Cerberus) with a robust Next.js/React dashboard and a mobile mission control (Pegasus), Aura provides a comprehensive solution for managing AI agents in a fully decentralized context.

\section{Prior Work and Foundation}
Aura is the practical realization of the theoretical frameworks established in prior research. Specifically, the integration of Hamiltonian State Space Duality (HSSD) within the Akasha 2 architecture \cite{meziani2026akasha} provides the low-latency, high-throughput foundation required for military-grade autonomous agent coordination.

\section{Architecture}
The Aura architecture is designed for high performance and low resource usage. The core components include:
\begin{itemize}
    \item \textbf{Aura Core}: The central orchestration layer.
    \item \textbf{Cerberus}: A Zig-based agent runtime (<1MB binary) for autonomous personas.
    \item \textbf{Pegasus}: An Android application for mobile mission control and Human-in-the-Loop (HITL) approvals.
    \item \textbf{Web Dashboard}: A Next.js/React interface for data visualization and management.
\end{itemize}

\section{Deployment Phases}

\subsection{Phase 1: Infrastructure}
The infrastructure layer is built on Debian 12 with a Python 3.11 environment. Key components include Ansible for control, Docker for containerization, and a robust database stack featuring PostgreSQL and Redis.

\subsection{Phase 2: Application Stack}
The application stack leverages modern technologies such as Kotlin for mobile development, Zig for the agent runtime, and Next.js for the web interface. We emphasize strict typing and modular architecture.

\subsection{Phase 3: Operational Workflows}
Aura supports complex operational workflows, including mission planning, task templates, and real-time status reporting via MQTT. Security is managed through OAuth 2.1 and XACML 3.1.

\section{UX Principles}
\begin{enumerate}
    \item \textbf{Operator First Design}: One-handed navigation, high-contrast UI, and voice command integration.
    \item \textbf{Resilience}: 30-day offline cache and redundant communication protocols.
    \item \textbf{Auditability}: Full mission logs and digital signatures for data provenance.
\end{enumerate}

\section{Conclusion}
Aura represents a significant step towards sovereign AI management. By providing a secure, auditable, and resilient platform, we enable researchers and developers to deploy AI agents with confidence.

\section*{Availability}
The source code for Aura is available under the MIT License on GitHub.

\begin{thebibliography}{9}
\bibitem{meziani2026akasha}
Meziani, Y. (2026). \textit{Akasha 2: Multimodal Architecture Integration with Hamiltonian State Space Duality}. arXiv preprint arXiv:2601.06212.
\end{thebibliography}

\end{document}
